\section{Introduction}\label{sec:introduction}

The 200-yard freestyle occupies an awkward and interesting place in
competitive swimming: long enough that energy management decides the outcome,
short enough that the anaerobic reserve is a first-order quantity rather than
a correction. Coaches prescribe pacing for it with strong conviction and
little shared formalism. This paper asks the question underneath those
prescriptions: \emph{what pacing strategy minimizes 200-yard freestyle race
time under realistic energetic and fatigue constraints, and how closely do
successful real-world performances match it?}

The mathematical study of race pacing begins with
\citet{keller1973theory,keller1974optimal}, who derived optimal running
strategies from a bounded propulsive force, linear resistance, and a finite
energy reserve with constant aerobic resupply. The line he opened was
extended with richer, hydraulic energetics \citep{behncke1993mathematical},
fitted to world records from 50~m to 275~km \citep{woodside1991optimal},
augmented with explicit fatigue states \citep{mathis1989fatigue}, and, in the
modern optimal-control treatment of \citet{aftalionbonnans2014}, shown to
favor deliberate velocity variation once anaerobic energy can be partially
re-created during deceleration. Swimming has been touched by this line
exactly once at the level of formal optimal control: \citet{maronski1996minimum}
derived a minimum-time swim profile --- glide, cruise, terminal kick --- from
a two-equation model. What the swimming literature has in abundance instead
is careful \emph{description}: elite 200-m freestyle finalists swim a fast
first lap followed by a nearly flat back half \citep{robertson2009analysis},
fast-start profiles dominate at every level, and pacing skill discriminates
finalists from non-qualifiers. Description and optimization have largely not
met, and where models exist they are typically fitted to records rather than
tested against the split sheets of ordinary competitive races.

This paper tries to make them meet, for one event, one course, and one
population: the 200-yard short-course freestyle swum by male competitive
swimmers aged 15--18. The contributions are four.

\paragraph{A theorem where a simulation is expected.} Under a fixed energy
budget and a velocity-only cost function, the four-split problem is a convex
program whose unique optimum is exactly even pacing, for every cost exponent
$p>1$ and every parameter choice (Section~\ref{sec:theorems}). Even pacing is
therefore not a finding about our parameters but a structural property of
budget-constrained racing, and every departure of real races from even pacing
is evidence about mechanisms \emph{outside} that structure.

\paragraph{Shape--engine separation.} With position-coupled economy decay the
optimal split allocation has a closed form, $t_i \propto w_i^{1/p}$, whose
striking consequence is that aerobic capacity, anaerobic capacity, drag and
efficiency have \emph{no effect at all} on the optimal shape
(Section~\ref{sec:theorems}). A swimmer with a bigger engine should swim the
same shape, faster. The popular intuition that stronger athletes have earned
a faster opening lap has no support in this model class.

\paragraph{Competing mechanisms as competing predictions.} We formalize five
fatigue mechanisms --- none of them exotic, all of them present somewhere in
the running literature --- and show they collapse onto three qualitatively
distinct split shapes: exactly even, strongly negative, and two
distinguishable kinds of positive (Section~\ref{sec:competing}). This turns
``which fatigue story is true'' into a sign-and-shape question that a few
dozen official split sheets can already inform.

\paragraph{Pre-registered confrontation with real races.} We built a
verified data pipeline for official age-group race results (anonymized at
ingest; swimmer-grouped train/test split; analysis plan and amendment log
frozen before any fitting) and report the confrontation honestly, including
the finding that currently dominates all others: the comparison is decided
less by fatigue physics than by the dive-start credit, whose
literature-anchored uncertainty band spans the distance between the
competing models (Sections~\ref{sec:data}--\ref{sec:results-empirical}).

Two boundaries are drawn deliberately. We optimize a clock, not a field:
tactical racing, drafting and psychology are outside the model. And we treat
within-lap velocity as constant, which excludes by assumption the
oscillation mechanisms of \citet{aftalionbonnans2014}; split-level data
cannot resolve within-lap dynamics in either direction, so this is a stated
limitation rather than a hidden one (Section~\ref{sec:limitations}).
